% RETHLAS-VERIFIED: P05.5 conditional model bridge; run 20260901T085240Z_45492fcf05c5_a5f577dc
% Exact target components: statement 8425157357c74519c30765e3e5f060688f71ae4a8be768f9516418aa28d357c9;
% proof 67180ff335763d08ec1eb9e3d2d72657a42002e8bea16b35b0f2147fc190a291.
\begin{proposition}
\label{prop:p055-model-calibration-transfer}
Assume the BFJ field and variety hypotheses of
\eqref{eq:ch24-standing}, with the variety denoted by $Y$ and
$\dim Y=d$.  Let $\mathbf X$ be a pure $d$-dimensional balanced rational
polyhedral space, let $\gamma$ be a piecewise-affine polyhedrally convex
reference function positive at infinity, and let
$a_1,\ldots,a_d\in\mathrm{PAPC}(\mathbf X,\gamma)$.  Write
\begin{equation}
 \mu_{\mathrm{poly}}
 :=\MA_{\mathrm{poly}}(a_1,\ldots,a_d)
 =\sum_{q\in Q}m_q\delta_q,
 \qquad m_q>0,
 \qquad \sum_{q\in Q}m_q=\deg(\gamma),
 \label{eq:p055-model-polyhedral-atoms}
\end{equation}
where $Q\subseteq\mathbf X$ is finite.

Let $j:\mathbf X\hookrightarrow\overline{\mathbf X}$ be an open embedding
into a locally compact Hausdorff space.  Let $L$ be an ample line bundle on
$Y$, let $\|\cdot\|_0$ be a semipositive model metric on $L$, and put
$\theta=c_1(L,\|\cdot\|_0)$.  For $1\leq i\leq d$, let
\begin{equation}
 \|\cdot\|_i=\mathrm e^{-u_i}\|\cdot\|_0
 \label{eq:p055-model-relative-metrics}
\end{equation}
be a semipositive model metric on $L$, and set
$\theta_i=\theta+\ddc u_i$.  Suppose that all the $\theta_i$ are determined
on one normal flat projective model $\mathcal Y$, write
\[
 \mathcal Y_s=\sum_E b_EE,
 \qquad x_E=b_E^{-1}\ord_E,
\]
and define the nonnegative intersection weights
\begin{equation}
 w_E:=b_E
 \bigl(\theta_{1,\mathcal Y}|_E\cdots
       \theta_{d,\mathcal Y}|_E\bigr)\geq0.
 \label{eq:p055-model-intersection-weights}
\end{equation}

Let $\overline h:Y^{\mathrm{an}}\to\overline{\mathbf X}$ be a proper
continuous map, put
\[
 U:=\overline h^{-1}(\mathbf X),
 \qquad h:=\overline h|_U,
\]
and assume that
\begin{equation}
 u_i|_U=(a_i-\gamma)\circ h
 \qquad(1\leq i\leq d).
 \label{eq:p055-model-lift}
\end{equation}
Assume moreover the exact mass identity
\begin{equation}
 (L^d)=\deg(\gamma)
 \label{eq:p055-model-mass-match}
\end{equation}
and the componentwise calibration
\begin{equation}
 \sum_{\substack{E:\ x_E\in U\\ h(x_E)=q}}w_E=m_q
 \qquad(q\in Q).
 \label{eq:p055-model-component-calibration}
\end{equation}

Then the BFJ mixed model measure, the Chambert--Loir mixed measure, and the
CLD product of the formal first Chern currents all coincide with
\begin{equation}
 \mu_{\mathrm{NA}}
 :=\MA_\theta(u_1,\ldots,u_d)
 =\sum_Ew_E\delta_{x_E}.
 \label{eq:p055-model-common-measures}
\end{equation}
They satisfy
\begin{equation}
 \mu_{\mathrm{NA}}(Y^{\mathrm{an}}\setminus U)=0
 \label{eq:p055-model-boundary-zero}
\end{equation}
and
\begin{equation}
 \overline h_*\mu_{\mathrm{NA}}
 =j_*\mu_{\mathrm{poly}}
 \label{eq:p055-model-global-comparison}
\end{equation}
as finite positive Radon measures on $\overline{\mathbf X}$.

The conclusion is conditional on the proper polarized realization, the
global model lifts \eqref{eq:p055-model-lift}, semipositivity, the mass
identity \eqref{eq:p055-model-mass-match}, and the intersection-theoretic
calibration \eqref{eq:p055-model-component-calibration}.  It asserts none
of those inputs for an arbitrary balanced polyhedral space.  It also makes
no assertion about a GK delta-form product on the original analytification
under these discretely valued field hypotheses.
\end{proposition}

\begin{proof}
The common-locus comparison
\Cref{thm:bfj-cld-common-locus-synthesis} identifies the BFJ mixed model
measure with the Chambert--Loir measure and with the product of the formal
CLD first Chern currents.  On the common
model $\mathcal Y$, their common value is the explicit model measure
\eqref{eq:model-mixed-measure}, namely
\[
 \sum_E b_E
 \bigl(\theta_{1,\mathcal Y}|_E\cdots
       \theta_{d,\mathcal Y}|_E\bigr)\delta_{x_E}
 =\sum_Ew_E\delta_{x_E}.
\]
This proves \eqref{eq:p055-model-common-measures}.  Relative nefness of the
models defining the semipositive metrics gives $w_E\geq0$.  By
\Cref{prop:model-ma-basic}, the total mass is the generic mixed
intersection number.  Since every metric is a metric on the same line
bundle $L$, this gives
\begin{equation}
 \sum_Ew_E=\mu_{\mathrm{NA}}(Y^{\mathrm{an}})=(L^d).
 \label{eq:p055-model-na-total-mass}
\end{equation}

For distinct points $q\in Q$, the sets of components occurring in the
sums in \eqref{eq:p055-model-component-calibration} are disjoint.  Summing
that calibration over $Q$, and then using
\eqref{eq:p055-model-polyhedral-atoms},
\eqref{eq:p055-model-mass-match}, and
\eqref{eq:p055-model-na-total-mass}, yields
\begin{align*}
 \sum_{q\in Q}
 \sum_{\substack{E:\ x_E\in U\\h(x_E)=q}}w_E
 &=\sum_{q\in Q}m_q
 =\deg(\gamma)
 =(L^d)
 =\sum_Ew_E.
\end{align*}
Every $w_E$ is nonnegative.  Hence each component omitted from the
double sum has weight zero.  In particular, if $x_E\notin U$, then
$w_E=0$, which proves \eqref{eq:p055-model-boundary-zero}.  The same
argument also kills every uncalibrated atom lying in $U$ whose image does
not belong to $Q$.

Let $B\subseteq\overline{\mathbf X}$ be Borel.  Using the atomic formula,
the preceding vanishing, and the componentwise calibration, we obtain
\begin{align*}
 (\overline h_*\mu_{\mathrm{NA}})(B)
 &=\sum_Ew_E\,\mathbf 1_B\bigl(\overline h(x_E)\bigr)\\
 &=\sum_{q\in Q}
   \left(
    \sum_{\substack{E:\ x_E\in U\\h(x_E)=q}}w_E
   \right)\mathbf 1_B(q)\\
 &=\sum_{q\in Q}m_q\mathbf 1_B(q)
  =\mu_{\mathrm{poly}}(B\cap\mathbf X)
  =(j_*\mu_{\mathrm{poly}})(B).
\end{align*}
This is \eqref{eq:p055-model-global-comparison}.  All measures are finite
positive Radon measures: this is part of the model-measure construction on
$Y^{\mathrm{an}}$, and the last equality identifies the extension by zero
with the proper pushforward.

The identities \eqref{eq:p055-model-lift} record that the formal metrics
are the proposed lifts of the polyhedral potentials.  The deduction from
the model measure to the pushforward identity uses their local geometric
content precisely through the independently stated calibration
\eqref{eq:p055-model-component-calibration}; it does not infer that
calibration from balancing alone.
\end{proof}
